
\subsubsection{Summary of Findings}

The existence proof (RQ1) demonstrates that task-category clustering in LoRA adapter geometry is statistically detectable under controlled conditions (Cohen's d = 0.765, p < $10^{-27})$. The correlation with FLAN taxonomy (RQ2, $\rho$ = 0.389) indicates that the geometric structure reflects semantic relationships between tasks.

The layer-type analysis yielded a negative result: all seven layer types show similar clustering strength (Cohen's d range 0.745--0.783), and no layer type exceeds the 0.8 threshold. Task-relevant geometric information appears to be distributed uniformly across the transformer architecture rather than concentrated in specific layers.

\subsubsection{Limitations}

\textbf{Training stochasticity.} The P3 control failure (ratio = 0.89) indicates that within-task variance across random seeds is comparable to within-category variance across tasks. This limits the practical utility of geometric signatures for fine-grained task discrimination. The geometric signatures may reflect category-level structure rather than precise task-level fingerprints.

\textbf{Single model family.} Experiments use only TinyLlama-1.1B. Generalization to other architectures and larger model scales is not established.

\textbf{Coarse taxonomy.} The binary FLAN taxonomy (same category vs. different category) provides only coarse similarity labels. The moderate correlation ($\rho$ = 0.389) may reflect this limitation.

\textbf{No baseline comparison.} This work validates the existence of geometric structure but does not compare Grassmann distance to alternative task similarity methods such as task embedding similarity or activation-based metrics.

\subsubsection{Scope Clarification}

The findings establish that geometric structure exists and correlates with task taxonomy. The results do not demonstrate that Grassmann distance is superior to or competitive with existing task similarity methods. Such comparison requires future work.
